\documentclass[a4paper,12pt]{article}
\usepackage[T2A]{fontenc}
\usepackage[utf8]{inputenc}
\usepackage{cmap}
\usepackage[english, russian]{babel}
\usepackage{wrapfig}
\usepackage[backend=biber,style=gost-numeric]{biblatex} % стиль ГОСТ
\addbibresource{literatura.bib}
\usepackage{misccorr} % в заголовках появляется точка, но при ссылке на них ее нет
\usepackage{amssymb,amsfonts,amsmath,amsthm}  
\usepackage{indentfirst}
\usepackage[usenames,dvipsnames]{color} 
\usepackage[unicode,hidelinks]{hyperref}
% \hypersetup{%
%     pdfborder = {0 0 0}
\usepackage{multirow}
\usepackage{makecell,multirow} 
\usepackage{ulem}
%\usepackage{hyperref}
%\usepackage{graphicx,wrapfig}
%\renewcommand\epsilon{\varepsilon}
% \renewcommand\epsilon{\varepsilon}
%\graphicspath{{img/}}
\usepackage{geometry}
\geometry{left=1cm,right=1cm,top=2cm,bottom=2cm,bindingoffset=0cm,headheight=15pt}
\usepackage{fancyhdr} 
\linespread{1.2} 
\frenchspacing 
\renewcommand{\labelenumii}{\theenumii)} 
% \usepackage{caption}
%%%%%%%%%%%%%%%%%%%%%%%%%%%%%%%%%%%%%%%%%%%%%%%%%%%%%%%%%%%%%%%%%%%%%%%%%%%%%%%
%%%%%%%%%%%%%%%%%%%%%%%%%%%%%%%%%%%%%%%%%%%%%%%%%%%%%%%%%%%%%%%%%%%%%%%%%%%%%%%

\def\labauthor{Сарафанов И.Г.}
\def\labauthors{\labauthor}
\def\labnumber{127}
\def\labtheme{Таблица производных и интегралов}

%%%%%%%%%%%%%%%%%%%%%%%%%%%%%%%%%%%%%%%%%%%%%%%%%%%%%%%%%%%%%%%%%%%%%%%%%%%%%%%
	%применим колонтитул к стилю страницы
\pagestyle{fancy} 
	%очистим "шапку" страницы
\fancyhead{\LaTeX} 
	%слева сверху на четных и справа на нечетных
\fancyhead[L]{\labauthors} 
	%справа сверху на четных и слева на нечетных
\fancyhead[R]{Яковлев \textnumero 190} 
	%очистим "подвал" страницы
\fancyfoot{} 
	% номер страницы в нижнем колинтуле в центре
\fancyfoot[C]{\thepage} 
%%%%%%%%%%%%%%%%%%%%%%%%%%%%%%%%%%%%%%%%%%%%%%%%%%%%%%%%%%%%%%%%%%%%%%%%%%%%%%%
\usepackage{float}
\usepackage[mode=buildnew]{standalone}
\usepackage{tikz} 
% \usepackage{subcaption}
\usepackage{tikz,csvsimple,physics}
\makeatother
%\usepackage{booktabs}
\usepackage{pgfplots, pgfplotstable}
\usepackage[outline]{contour}
\usepackage{tocloft}
\renewcommand{\cftsecleader}{\cftdotfill{\cftdotsep}}


% вот этот удивительный мир

\begin{document}
\begin{center}
    \textbf{Домашняя работа \textnumero 1}
\end{center}

\textit{Дано: $h = 19{,}6$ м., $S_1=1000$ м., найти $S_2$.}
\begin{figure}[h]
    \centering
    \begin{tikzpicture}
    \draw[->] (0,0) -- (10,0) node[anchor=north] {$x$};
    \draw[->] (0,0) -- (0,5) node[anchor=east] {$y$};
    %\draw[thick] (0,4) node[left] {$A$} -- (0,0) node[below] {$O$} -- (4,0) node[below] {$B$};
    %\draw[dashed] (4,0) -- (4,4) node[right] {$C$};
    %\node at (2,-0.5) {1000 м};
    %\node at (-0.5,2) {19.6 м};
\end{tikzpicture}
\end{figure}

Из законов кинематики свободного падения получаем:
\[
h = \frac{g t^2}{2} \, \Rightarrow \, 19{,}6 = \frac{9{,}8 t^2}{2}
\, \Rightarrow \,
t = 2 \text{ c}.
\]

Горизонтальная скорость снаряда непосредственно перед взрывом:
\[
v_x = \frac{1000}{2} = 500 \text{ м/с}.
\]

Вертикальная скорость первой отколовшейся части:
\[
19{,}6 = - v_1 - 4{,}9
\Rightarrow
v_1 = 14{,}7 \text{ м/с}.
\]

По закону \textbf{сохранения импульса:}
\[
m v_{1y} + m v_{2y} = 0
\Rightarrow
v_{2y} = 14{,}7 \text{ м/с}.
\]

По горизонтали:
\[
m \cdot 0 + m v_{2x} = 2m \cdot 500
\Rightarrow
v_{2x} = 1000 \text{ м/с}.
\]

Из кинематики находим время движения второй части после взрыва:
\[
t_2 = 4 \text{ c}.
\]

Горизонтальное перемещение после взрыва:
\[
S = 1000 \cdot 4 = 4000 \text{ м}.
\]

Полное расстояние:
\[
S_2 = 1000 + 4000 = 5000 \text{ м}.
\]

\[
\boxed{S_2 = 5000 \text{ м}}
\]

\end{document}